\documentclass[11pt,a4paper]{article}

\usepackage[utf8]{inputenc}
\usepackage[T1]{fontenc}
\usepackage{lmodern}
\usepackage[margin=22mm,top=24mm,bottom=24mm]{geometry}
\usepackage{booktabs}
\usepackage{array}
\usepackage{longtable}
\usepackage{xcolor}
\usepackage{colortbl}
\usepackage{enumitem}
\usepackage{microtype}
\usepackage{seqsplit}
\setlength{\emergencystretch}{4em}
\sloppy
\hbadness=10000
\hfuzz=4pt
% Wrap long filenames with breakable points
\newcommand{\fpath}[1]{\texttt{\seqsplit{#1}}}
\usepackage{titlesec}
\usepackage{fancyhdr}
\usepackage{parskip}
\usepackage[hidelinks,colorlinks=true,linkcolor=sambpblue,urlcolor=sambpblue]{hyperref}

% --- Colour palette mirroring the SAMBPS DTaaS execution manual ----------
\definecolor{sambpblue}{HTML}{1F3D7A}
\definecolor{sambplight}{HTML}{D9E2F3}
\definecolor{sambpgrey}{HTML}{4A4A4A}
\definecolor{sambpaccent}{HTML}{F39200}

% --- Headings ---------------------------------------------------------------
\titleformat{\section}
  {\color{sambpblue}\Large\bfseries}{\thesection.}{0.6em}{}
\titleformat{\subsection}
  {\color{sambpblue}\large\bfseries}{\thesubsection}{0.6em}{}
\titleformat{\subsubsection}
  {\color{sambpgrey}\normalsize\bfseries}{\thesubsubsection}{0.6em}{}

% --- Headers / footers ------------------------------------------------------
\pagestyle{fancy}
\fancyhf{}
\renewcommand{\headrulewidth}{0.4pt}
\renewcommand{\footrulewidth}{0.4pt}
\fancyhead[L]{\footnotesize\color{sambpblue}\textbf{SAMBPS DTaaS}\ \textendash\ Fault Location Identification}
\fancyhead[R]{\footnotesize\color{sambpgrey}Project Execution Plan v1.0}
\fancyfoot[L]{\footnotesize Anoop Eluvathingal -- Chief Scientist \& CEO, SAMBPS DTaaS}
\fancyfoot[R]{\footnotesize\thepage}

% --- Custom column types ----------------------------------------------------
\newcolumntype{L}[1]{>{\raggedright\arraybackslash}p{#1}}
\newcolumntype{C}[1]{>{\centering\arraybackslash}p{#1}}

\setlist[itemize]{leftmargin=*,topsep=2pt,itemsep=1pt}
\setlist[enumerate]{leftmargin=*,topsep=2pt,itemsep=1pt}

\begin{document}

% ===========================================================================
% Cover
% ===========================================================================
\begin{titlepage}
\thispagestyle{empty}
\noindent\colorbox{sambpblue}{\parbox{\dimexpr\textwidth-2\fboxsep\relax}{%
  \vspace{0.6cm}
  \color{white}\Large\textbf{SAMBPS DTaaS}\\[2pt]
  \color{sambplight}\normalsize Smart Adaptive Model-Based Protection \& Control Systems\\[1.2cm]
  \color{white}\Huge\textbf{Fault Location Identification}\\[6pt]
  \color{sambplight}\Large Project Execution Plan v1.0\\[6pt]
  \color{white}\large HIF transfer-function locator -- 48-week, 252-PD operating plan
  \vspace{0.4cm}
}}

\vspace{1.6cm}

\noindent{\large\textbf{Subject under execution}}\\[3pt]
\noindent
Single-ended joint estimation of HIF location $\alpha$ and arc resistance $R_x$
via power-frequency admittance identification with dual-channel noise modelling
(IEEE Access manuscript, Arjundas K.\ \& K.\ Shanthi Swarup, 2026).

\vspace{0.7cm}

\noindent{\large\textbf{Reference baseline}}\\[3pt]
\noindent
\texttt{HIF\_TF\_Execution\_Manual.pdf} (v3) -- 18\,pp. internal SAMBPS execution
manual that operationalises the v2 deep-research project plan into day-level
WBS, RACI, acceptance specs, schedule and repository inventory.  This document
mirrors the manual's structure and is the canonical project-execution plan for
the new \texttt{04\_code/sambp/fault\_location\_id/} sub-project, sized and
named consistently with the four existing SAMBP sub-projects
(\texttt{sync\_oc}, \texttt{transformer\_diff}, \texttt{line\_diff},
\texttt{bus\_diff}).

\vspace{0.7cm}

\noindent{\large\textbf{Posture}}\\[3pt]
\noindent
Phase 0 unlocks IEEE Access camera-ready in $\approx 14$ person-days over weeks
W0--W2.  Phases 1--5 retire the two structural risks identified by the v2 plan
(corrected CRLB, continuously parametrised optimiser) and convert the
manuscript into a literature-anchored, three-phase, hardware-tested SAMBPS
DTaaS Protection-Validation module with a second journal paper and an
institutional HIL pilot.  Total horizon 48 weeks; total effort
$\approx 252$\,person-days; staffing one full-time lead engineer + one
half-time second engineer + advisor time per the RACI in \S\,5.

\vfill
\noindent\colorbox{sambplight}{\parbox{\dimexpr\textwidth-2\fboxsep\relax}{%
  \vspace{0.15cm}
  \color{sambpblue}\textbf{Prepared by} Anoop Eluvathingal -- Chief Scientist \& CEO, SAMBPS DTaaS\\[1pt]
  \color{sambpgrey}\small Issued: 09 May 2026 \quad\textbar\quad Confidential -- internal execution plan\\[1pt]
  \small Located in: \texttt{04\_code/sambp/fault\_location\_id/docs/}
  \vspace{0.15cm}
}}
\end{titlepage}

\tableofcontents
\newpage

% ===========================================================================
\section{Purpose, scope and how to use this plan}
% ===========================================================================

This Execution Plan is the day-level operating document for the new
SAMBPS DTaaS sub-project \texttt{fault\_location\_id} (Fault Location
Identification).  It mirrors the structure, notation and decision-gate
discipline of the SAMBPS HIF-TF Execution Manual v3 and is written so the
lead engineer (Arjundas K.) can pick it up on Monday morning, drive
Phase~0 to camera-ready, and continue without ambiguity through Phases
1--5 to the second journal paper and the SAMBPS DTaaS Protection-Validation
module v1.0.

\paragraph{How to read this plan}

\begin{itemize}
  \item \textbf{\S\,2 Executive summary} -- five-phase posture, critical work
    packages, decision gates.
  \item \textbf{\S\,3 Section-by-section audit} -- nine manuscript-level
    shortcomings, each tagged to the Work Package(s) and Deliverable(s)
    that close it.
  \item \textbf{\S\,4 Work Breakdown Structure} -- every WP from the v2/v3
    plan decomposed into numbered sub-tasks with effort estimates in
    person-days (PD) and dependencies.
  \item \textbf{\S\,5 RACI matrix} -- nine-actor responsibility map across
    SAMBPS, IIT-Madras, NUS, NTU, TU Dortmund, Amprion and the IED
    industry partners.
  \item \textbf{\S\,6 Acceptance test specifications} -- the eight
    deliverables D-A through D-H locked to quantitative pass/fail
    predicates that CI executes on every commit.
  \item \textbf{\S\,7 Schedule} -- W0--W2 day-by-day; W2--W48 week-by-week.
  \item \textbf{\S\,8 Communication plan} -- daily / weekly / phase-end
    cadence and the decision-gate review template.
  \item \textbf{\S\,9 Repository inventory} -- the
    \texttt{fault\_location\_id} folder skeleton (created alongside this
    plan) and the public artefact map.
  \item \textbf{\S\,10 Risk register} -- ten-row register imported from the
    v2 plan and tightened to this sub-project.
  \item \textbf{\S\,11 KPI dashboard} -- 17 KPIs from the v2 plan with
    current values and SAMBPS DTaaS v1.0 targets.
  \item \textbf{\S\,12 Bottom line and approval block.}
\end{itemize}

\paragraph{Notation.}
\textbf{WPx.y} = work package from the v2 plan
(e.g.\ WP1.6 = re-derive the FIM via the proper-complex-Gaussian-ratio
framework).  \textbf{D-A \ldots\ D-H} = v2 deliverables.
\textbf{D0 \ldots\ D5} = decision gates at the end of each phase.
\textbf{R1 \ldots\ R10} = v2 risk register rows.  \textbf{O1 \ldots\ O9} =
v2 objectives.  Manual references like ``v3 \S\,3.7'' point to subsections
of \texttt{HIF\_TF\_Execution\_Manual.pdf}.

% ===========================================================================
\section{Executive summary -- execution posture}
% ===========================================================================

\renewcommand{\arraystretch}{1.18}
\begin{longtable}{L{1.4cm}L{4.0cm}L{5.6cm}L{4.0cm}}
\toprule
\rowcolor{sambplight}
\textbf{Phase} & \textbf{Posture} & \textbf{Critical work packages} & \textbf{Decision gate}\\
\midrule
\endhead

\textbf{Phase 0}\newline W0--2 &
Editorial + reproducibility hardening; published as soon as Phase 0 closes.
All work parallelisable across two engineers. &
WP0.1 metric harmonisation; WP0.2 prior-art restructure; WP0.3 references
+ DOI pass; WP0.4 public repo + CI; WP0.5 Appendix A
($\pi$-derivation + $\partial H/\partial\theta$). &
\textbf{D0:} IEEE Access camera-ready v2.0 frozen; CI green.\\

\textbf{Phase 1}\newline W2--10 &
Two work-streams in parallel: (a)~cross-platform validation;
(b)~corrected CRLB. &
WP1.1--1.4 PSCAD + EMTP + 50-section reference + 720-rerun;
WP1.5 Monte-Carlo; WP1.6 proper-complex-Gaussian-ratio FIM. &
\textbf{D1:} Mean location error $<2$\,\% across all simulators at
$\mathrm{SNR}_I\geq 30$\,dB; corrected CRLB visualised.\\

\textbf{Phase 2}\newline W6--16 (overlaps Phase 1) &
Mathematical work to retire the 39.44\,\% section-modelling-error
ceiling; gated on Phase 1 D1. &
WP2.1 closed-form $H$; WP2.2 $\partial H/\partial\alpha$ analytic;
WP2.3 verify vs 50-section; WP2.4 analytical gradients in optimiser;
WP2.5 re-run 720 + Monte-Carlo. &
\textbf{D2:} Modelling error vs 50-section $<5$\,\%;
estimator improvement $\geq 30$\,\% at noisy SNR.\\

\textbf{Phase 3}\newline W10--24 &
Generalisation to three-phase, branched, multi-section, IEEE feeders;
replace single-bin DFT. &
WP3.1--3.4 three-phase $Y_{abc}$ + IEEE 13/34/123 + fault types;
WP3.5 Taylor-Fourier; WP3.6 multi-port FIM; WP3.7 CNRS dataset. &
\textbf{D3:} Mean location error $<3$\,\% on IEEE 34-node at SNR
$\geq 30$\,dB.\\

\textbf{Phase 4}\newline W18--30 (overlaps Phase 3) &
Field-grade impairments + numerical competitor benchmark + arc-class
diversification. &
WP4.1 impairments; WP4.2--4.4 Cassie-Mayr-Kizilcay + Wang-2020 +
Torres-2022; WP4.5 four-competitor benchmark. &
\textbf{D4:} $<5$\,\% location error at SNR $\geq 30$\,dB across all
impairment classes; numerical superiority over $\geq 2$ of 4
competitors.\\

\textbf{Phase 5}\newline W26--48 &
HIL pilot + DTaaS productisation; institutional partner engagement
throughout. &
WP5.1 HIL platform; WP5.2 real IED + IEC 61850-9-2 SV;
WP5.3 live HIF arc; WP5.4 DTaaS module; WP5.5 second journal paper. &
\textbf{D5:} Real-IED end-to-end estimate latency $<5$ cycles;
institutional pilot signoff; SAMBPS DTaaS v1.0 release; second journal
submitted.\\

\bottomrule
\end{longtable}

% ===========================================================================
\section{Section-by-section manuscript audit}
% ===========================================================================

The audit below mirrors the manuscript's running order.  Each row
identifies what is currently written, the shortcoming after comparison
with the v2 plan, and the upgrade actions tagged with the v2 work
package(s) and deliverable(s) that close them.

\renewcommand{\arraystretch}{1.18}
\begin{longtable}{L{4.4cm}L{4.6cm}L{6.0cm}}
\toprule
\rowcolor{sambplight}
\textbf{Existing writing} & \textbf{Shortcoming vs v2 plan} & \textbf{Upgrade action (WP / deliverable)}\\
\midrule
\endhead

\textbf{Title.}
``HIF Location and Resistance Estimation \ldots Transfer Function
Identification at 50\,Hz with Dual-Channel Noise Modelling.'' Single-ended TF
identification.
&
Title leads with the technique, not the deliverable; ``at 50 Hz'' reads as
a constraint; the genuine novelty (joint $\alpha + R_x$) is buried.
&
\textbf{WP0.1.} Title to: ``Single-Ended Joint Estimation of HIF Location
and Arc Resistance via Power-Frequency Admittance Identification with
Dual-Channel Noise Modelling.'' Add training-free / single-frequency to
keywords.  Closes \textbf{D-A}.\\

\textbf{Abstract.}
Method, two-stage optimiser, 720-case AWGN sweep, headline 0.009\,\% /
1.18\,\%, 39.44\,\% vs 50-section, CRLB analysis.
&
Conflates location\,\% vs modelling\,\%; CRLB headline
(``within 15\,\% of CRLB at SNR $\geq 40$\,dB'') absent;
$R_x$ error envelope absent; no motivation paragraph; sampling
config not stated.
&
\textbf{WP0.1.} Rewrite to $\leq 250$ words: open with PSRC D15 + NREL 2023
motivation; disambiguate modelling vs estimation error; add CRLB headline;
add $R_x$ envelope; state $f_s$ / $N_s$ / window.  Closes \textbf{D-A}.\\

\textbf{Introduction.}
Four-category taxonomy; single-ended single-frequency operation positioned
as the contribution; five contributions listed.
&
Taxonomy is dated; should be 7 families per Wang--Lopes-2023 EPSR
($\mu$-PMU, TW, morphology / wavelet, GAN / GNN, statistical missing).
No real-world fatality evidence; no CPUC SB 901 / PSPS framing.
&
\textbf{WP0.1 + WP0.2.} Restructure prior-art into 7 families with tabular
synthesis; add 1-paragraph wildfire / safety motivation; add
forward-reference roadmap.  Closes \textbf{D-A}.\\

\textbf{Distribution-line parameterisation (II.A).}
11\,kV, 100\,km feeder split at $\alpha\in(0,1)$.  $R'=0.0728$\,$\Omega$/km,
$L'=4R'$, $C'=3R'$.  Single-section, single-feeder, single-phase.
&
$L'=4R'$ and $C'=3R'$ are dimensionally heterodox; Saha-2010 and
Wang--Lopes-2023 require explicit units.  Single-section / single-phase
neglects the branched / DG-rich reality covered by Phase 3.
&
\textbf{WP0.5.} Add explicit numeric values
($L'=0.927$\,mH/km, $C'=11.6$\,nF/km) with citation; dimensional-check
sentence; scope statement on single-feeder restriction.
\textbf{WP3.1--3.3.} Generalise to three-phase $Y_{abc}$, IEEE 13/34/123.
Closes \textbf{D-A, D-D}.\\

\textbf{Two-section $\pi$-model state-space (II.B).}
$x = [V_{C1}, V_{C2}, I_{L1}, I_{L2}]^T$; 4$\times$4 $A$ matrix printed in
Eq.~(3); $A_{11} = -1/(R_x C_1)$ continuously differentiable in
$(\alpha, R_x)$.
&
Derivation correct but not appended; $\partial H/\partial\alpha$ and
$\partial H/\partial R_x$ not written out, although both the gradient
solver and the FIM need them.
&
\textbf{WP0.5.} Appendix A with annotated $\pi$-circuit, full KVL/KCL
$\rightarrow$ state-space derivation, symbolic
$\partial H/\partial\alpha$ and $\partial H/\partial R_x$.  Verify in
MATLAB \texttt{sym/eig} regression test (R7 mitigation).
Closes \textbf{D-A}.\\

\textbf{Anti-parallel diode arc model (II.C).}
$V_{kp}=50$\,V, $V_{kn}=45$\,V, $R_{sp}=5\,\Omega$, $R_{sn}=6\,\Omega$,
$R_{\mathrm{off}}=10^{6}\,\Omega$, $\varepsilon=10^{-3}$\,A;
quasi-static $R_x$ over one cycle.
&
Parameter provenance gap (Emanuel-1990 sandy-soil values are kV-range
$V_p / V_n$, tens-to-thousands $\Omega$ $R_p / R_n$;
Santos-2022 surface-resolved table contradicts the values); single arc
class; random reignition not modelled.
&
\textbf{WP0.5.} Cite the source of every parameter; if internal
calibration, document. \textbf{WP4.2--4.4.} Add Cassie--Mayr--Kizilcay,
Wang-2020 distortion-controllable, Torres-2022 stochastic-configurable;
report $\Delta$-error.  Closes \textbf{D-A, D-E}.  Mitigates R10.\\

\textbf{Section-model accuracy study (II.D).}
$N_s=2,5,10,20$ vs 50-section reference; 2-section mean 39.44\,\%, max
89.78\,\%; 10-section creates fault-node discretisation that breaks the
gradient solver; 2-section retained as deliberate trade-off.
&
\emph{Single most important residual issue.} Lopes-2023, Trew-2023,
Suonan-2023, Kang-2021 publish closed-form continuously-parametrised
distributed-parameter formulations in which the fault is inserted as a
shunt at $\alpha$ with two cascaded ABCD / Bergeron / hyperbolic blocks.
&
\textbf{WP2.1--2.5.} Derive $H(j\omega_0; \alpha, R_x)$ from cascaded
ABCD/Bergeron; prove $\partial H/\partial\alpha$ exists analytically;
reproduce 50-section reference within 1\,\%; replace FD with analytical
gradients in optimiser; re-run 720 + Monte-Carlo.  Closes \textbf{D-C}.
Decision gate \textbf{D2}.\\

\textbf{Transfer-function identification (III).}
$H_{\mathrm{model}} = C(j\omega_0 I - A)^{-1}B$; DFT at 10\,kHz / 200
samples / $k_{50}=2$; $J=|\Delta H|^2 = [\Delta\mathrm{Re}]^2 +
[\Delta\mathrm{Im}]^2$.
&
Single-bin DFT biased under non-stationary HIF; single complex observation
just-determined for 2 reals but locally degenerate where
$\partial H/\partial\alpha$ and $\partial H/\partial R_x$ are colinear;
Euclidean cost weighting sub-optimal -- Kay-1993 prescribes $\Sigma^{-1}$.
&
\textbf{WP3.5.} Replace single-bin DFT with first-order Taylor-Fourier;
map $J(\alpha, R_x)$ on representative case; apply Hermann-Krener ORC.
\textbf{WP3.6 + WP1.6.} Adopt ML / whitened-Mahalanobis cost
$J_{\mathrm{ML}} = \Delta H^H \Sigma^{-1}\Delta H$.  Closes \textbf{D-D}.\\

\textbf{Two-stage optimiser (IV).}
Stage 1: $100\times50$ grid + top-3 multi-start + 8-neighbour steepest
descent.  Stage 2: gradient descent with central FD + Armijo line-search;
box constraints; 2000-iter cap.
&
Top-3 multi-start has no global-optimum guarantee; FD wasteful when
analytic $\partial H/\partial\alpha$ available;
no runtime / FLOP / wall-clock report; hyperparameters lack sensitivity
analysis; ``2-cycle response'' claim not validated.
&
\textbf{WP2.4.} Replace FD with analytical gradients.
\textbf{WP0.4 sub-tasks 6/7/8.} Global-optimum capture statistic over
720-grid; hyperparameter sensitivity sub-table; median + 95th-percentile
CPU time (\texttt{tic/toc} over 1000 calls).
Closes \textbf{D-A, D-C}.\\

\textbf{Simulation setup and noise (V).}
Full-factorial 720-case ($10\,\alpha\times 5\,R_x\times 4\,\mathrm{SNR}_V
\times 4\,\mathrm{SNR}_I$); independent dual-channel AWGN; \texttt{rng(42)};
IEEE HIF range 100--5000\,$\Omega$.
&
Self-consistent simulation framework (same model generates and is fitted);
AWGN-only; no CT/PT saturation; no off-nominal frequency; means only.
&
\textbf{WP1.1--1.4} re-run on PSCAD / EMTP / 50-section reference.
\textbf{WP1.5} 100 Monte-Carlo trials per cell.
\textbf{WP4.1} add five impairment classes.
\textbf{WP3.7} validate on CNRS-2024 IEEE-34 dataset.
Closes \textbf{D-B, D-D, D-E}.  Decision gates \textbf{D1, D3, D4}.\\

\textbf{Results and discussion (VI).}
Six panels: noiseless baseline; $\mathrm{SNR}_V$ sweep;
$\mathrm{SNR}_I$ sweep (dominant); $R_x$ error sweeps; estimated-vs-true
scatter; $\mathrm{SNR}_V \times \mathrm{SNR}_I$ and
$\alpha \times R_x$ heatmaps.
&
Figure captions truncated; heatmaps don't overlay CRLB envelope;
``systematic bias-free'' claim not numerically tested; no comparison
plot vs competitors.
&
\textbf{WP0.4 sub-task 4.} Patch every figure caption; add axis labels
with units. \textbf{WP1.6.} Overlay corrected CRLB envelope on Figs 2--5
and 8--9. \textbf{WP1.5.} Add per-cell numerical bias and 95\,\% CI.
\textbf{WP4.5.} Add comparison plot vs Paramo / Iurinic / Cui-Weng / Zeng.
Closes \textbf{D-B, D-E}.\\

\textbf{Comparison with existing methods (VII).}
Table 3 contrasts proposed TF method with seven prior-art classes;
categorical noise level Low/Med/High; comm requirement; location
capability.
&
No numerical head-to-head; Iurinic-2018 spectral, Zeng-2021 double-ended
HIF, Cui-Weng-2020 $\mu$-PMU, Lopes-2023 distributed-parameter all
missing or shallowly cited; categorical noise non-quantitative.
&
\textbf{WP4.5.} Re-implement and re-run four competitors on identical
720-case waveforms.  Publish numerical Table 3-bis: method, location\,\%
error, $R_x$\,\% error, CPU time, comm infrastructure, training-data
requirement, SNR floor for $<5$\,\% error.  Closes \textbf{D-E}.
Mitigates R6.\\

\textbf{Cram\'er-Rao Lower Bound (VIII).}
FIM under circular complex Gaussian noise on $H_{\mathrm{meas}}$;
four headline findings; near-source $\alpha < 0.2$ floor identified
as fundamental.
&
\emph{Most consequential audit finding.} The data noise model in V-B is
dual-channel AWGN on $V$ and $I$ ratioed to form $H_{\mathrm{meas}}$;
the resulting density on $H$ is non-Gaussian; the Gaussian-on-$H$ FIM
under-estimates the true bound and is valid only when
$|I|\gg \sigma_I$ -- the opposite of the HIF regime.
&
\textbf{WP1.6.} Re-derive the FIM via the proper-complex-Gaussian-ratio
framework (Kuruoglu-2018 ASCE) or a joint dual-channel FIM in $(V, I)$
space (Nehorai-Hawkes-2000; Kay-1993 Vol.~I Ch.~15) projected onto
$(\alpha, R_x)$.  Publish $\partial H/\partial\alpha,
\partial H/\partial R_x$ in Appendix A (WP0.5).  Add CRLB-vs-empirical
overlay plots.  Add Geary-Hinkley validity test.  Closes \textbf{D-B}.
Decision gate \textbf{D1}.  Mitigates R9.\\

\textbf{Conclusion and future work (IX).}
Summarises method and headline numbers; identifies two-section limitation
as principal residual; lists future work: continuously-parametrised model,
three-phase, HIL, non-Gaussian noise.
&
Future-work list is broad but un-prioritised, un-scheduled, un-anchored;
$R_x$-error envelope buried; no TRL framing; SAMBPS DTaaS productisation
absent.
&
\textbf{WP0.1.} Convert future-work into numbered roadmap with horizons
(12 / 24 / 36 mo) aligning to v2 phases; tighten $R_x$ envelope
statement; add 2--3 line outlook on DTaaS Protection-Validation module
integration; cite $\geq 3$ literature anchors per future-work item.
Closes \textbf{D-A}.\\

\textbf{References, figures, supplementary.}
20 references (one duplicate); 9 figures (caption truncation issues);
no public repo; no Data Availability Statement.
&
Thin reference set for an active literature; v2 bibliography offers
$\geq 80$ candidate entries; Wei-2020 PSCAD HIAF GitHub model and
CNRS-2024 IEEE-34 dataset exemplars.
&
\textbf{WP0.3.} De-duplicate Refs [2]/[10]; expand to $\approx 35$--$45$
entries grouped by family; DOI + IEEE-style pass; verify every reference.
\textbf{WP0.4.} Public GitHub/Zenodo repo with .m source, 720-case MAT,
figure scripts, symbolic $\partial H/\partial\theta$ derivation, README,
LICENSE, CI; add Data Availability Statement.  Closes \textbf{D-A}.\\

\bottomrule
\end{longtable}

% ===========================================================================
\section{Work Breakdown Structure (day-level decomposition)}
% ===========================================================================

Every WP is decomposed into numbered sub-tasks with effort estimates in
person-days (PD).  Effort assumes one experienced MATLAB / Python power-systems
engineer; halve elapsed time when two engineers work in parallel.

\subsection{Phase 0 -- Editorial + reproducibility (W0--2; total $\approx 14$ PD)}

\begin{longtable}{L{1.2cm}L{9.4cm}C{1.2cm}L{3.2cm}}
\toprule\rowcolor{sambplight}
\textbf{WP} & \textbf{Sub-task} & \textbf{Effort} & \textbf{Output}\\\midrule\endhead
WP0.1 & Title rewrite (\S\,3.1); abstract rewrite (\S\,3.2); conclusion roadmap
       (\S\,3.14); metric harmonisation across abstract / Sect.~VI / Sect.~IX. &
       1.5\,PD & frozen \fpath{HIF-TF\_v2.0\_manuscript.docx} + \texttt{changelog.md}\\
WP0.2 & Restructure prior-art into 7-family taxonomy; wildfire-motivation paragraph
       (NREL 2023, PSRC D15, Camp Fire, Black Saturday, CPUC SB 901);
       forward-reference roadmap. &
       1.0\,PD & revised Sect.~I\\
WP0.3 & Reference set expansion to $\approx 35$--$45$ entries; de-duplicate
       Refs [2]/[10]; DOI + IEEE-style pass; verify every entry against IEEE
       Xplore / Crossref. &
       1.5\,PD & \fpath{references.bib}\\
WP0.4 & Stand up public GitHub/Zenodo repo: \texttt{README}, \texttt{LICENSE} (MIT),
       \texttt{.gitignore}, GH Actions CI workflow (unit tests + symbolic
       \texttt{sym/eig} regression); push existing .m source + 720-case MAT +
       figure scripts.  Sub-tasks 4--8: figure caption patch; CRLB-overlay
       placeholder; global-optimum capture statistic; hyperparameter sensitivity
       table; \texttt{tic/toc} CPU report. &
       4.0\,PD & repository (private $\rightarrow$ public at D0); Zenodo DOI\\
WP0.5 & Appendix A: annotated $\pi$-circuit; full KVL/KCL $\rightarrow$ state-space
       derivation; symbolic $\partial H/\partial\alpha$ and
       $\partial H/\partial R_x$; arc-parameter provenance citations; explicit
       numeric line parameters with units; dimensional-check sentence. &
       3.0\,PD & \fpath{AppendixA\_derivation.tex} + \fpath{derive\_partials.m}\\
WP0.6 & IEEE Access camera-ready + supplementary upload; co-author signoff;
       submit. &
       1.0\,PD & submission ID; D-A achieved\\
\midrule
\multicolumn{2}{r}{\textbf{Phase 0 total}} & \textbf{14.0\,PD} & \\
\bottomrule
\end{longtable}

\subsection{Phase 1 -- Cross-platform + corrected CRLB (W2--10; $\approx 35$\,PD)}

\begin{longtable}{L{1.2cm}L{9.4cm}C{1.2cm}L{3.2cm}}
\toprule\rowcolor{sambplight}
\textbf{WP} & \textbf{Sub-task} & \textbf{Effort} & \textbf{Output}\\\midrule\endhead
WP1.1 & Build PSCAD model: 11\,kV / 100\,km feeder + anti-parallel diode arc +
       dual-channel CT/PT; export V/I waveforms at 10\,kHz to \texttt{.mat}. &
       6.0\,PD & \fpath{pscad/HIFL\_11kV\_100km.pscx}\\
WP1.2 & EMTP-RV mirror of WP1.1 (independent implementation by E2; or
       third-party EMTP user reviews PSCAD case). &
       5.0\,PD & \fpath{emtp/HIFL\_11kV\_100km.ecf}\\
WP1.3 & Pure-MATLAB 50-section reference state-space (data-only; not used by
       optimiser); generate waveforms across all 720 grid points. &
       2.0\,PD & \fpath{matlab/ref\_50section\_dataset.mat}\\
WP1.4 & Re-run optimiser (unchanged) on each independent dataset; compute mean
       / max location and $R_x$ error; quantify $\Delta$-error attributable to
       optimiser-vs-data model mismatch. &
       3.0\,PD & \fpath{results/phase1\_crossplatform.csv} + report\\
WP1.5 & 100-trial Monte-Carlo per grid cell across each dataset; mean /
       95th-percentile; per-cell numerical bias and 95\,\% CI. &
       4.0\,PD & \fpath{phase1\_montecarlo.mat} + per-cell CI table\\
WP1.6 & Re-derive FIM via proper-complex-Gaussian-ratio framework
       (Kuruoglu-2018) AND joint dual-channel FIM in $(V,I)$ projected onto
       $(\alpha, R_x)$ (Nehorai-Hawkes-2000); cross-check; Geary-Hinkley
       validity test; CRLB-vs-empirical overlay plots. &
       8.0\,PD & \fpath{crlb\_proper.m}, \fpath{crlb\_dualchannel.m},
       \fpath{fig/crlb\_overlay\_*.pdf}, \fpath{AppendixB\_correctedCRLB.tex}\\
WP1.7 & Phase 1 milestone document; arXiv preprint of cross-validation +
       corrected CRLB; D1 review. &
       3.0\,PD & arXiv preprint; D-B achieved\\
\midrule
\multicolumn{2}{r}{\textbf{Phase 1 total}} & \textbf{31.0\,PD} & rounded $\approx 35$ with contingency\\
\bottomrule
\end{longtable}

\subsection{Phase 2 -- Continuously parametrised optimiser (W6--16; $\approx 28$\,PD)}

\begin{longtable}{L{1.2cm}L{9.4cm}C{1.2cm}L{3.2cm}}
\toprule\rowcolor{sambplight}
\textbf{WP} & \textbf{Sub-task} & \textbf{Effort} & \textbf{Output}\\\midrule\endhead
WP2.1 & Derive cascaded ABCD / Bergeron / hyperbolic closed-form admittance
       $H(j\omega_0; \alpha, R_x)$ where the fault is a shunt at distance
       $\alpha$ between two distributed-parameter line sections (Lopes-2023;
       Trew-2023; Kang-2021; Pozar Microwave Eng.\ ABCD). &
       6.0\,PD & \fpath{H\_closedform.m} + \fpath{derive\_H\_closedform.m}\\
WP2.2 & Derive $\partial H/\partial\alpha$ and $\partial H/\partial R_x$ in
       closed form via analytic derivatives of $\cosh(\gamma\alpha\ell)$,
       $\sinh(\gamma\alpha\ell)$; MATLAB \texttt{sym/diff} verification. &
       3.0\,PD & \fpath{dH\_dalpha.m}, \fpath{dH\_dRx.m}\\
WP2.3 & Reproduce 50-section reference within 1\,\% across all $(\alpha, R_x)$
       grid; if $>1$\,\%, root-cause and iterate. &
       4.0\,PD & \fpath{results/phase2\_modelfit.csv}\\
WP2.4 & Replace FD with analytical gradients in optimiser; hyperparameter
       sensitivity sub-table; re-verify Armijo line-search behaviour. &
       4.0\,PD & updated \texttt{optimiser.m}; sensitivity table\\
WP2.5 & Re-run all 720 cases + Monte-Carlo on the new optimiser; compare to
       Phase 1 baseline; report estimator improvement. &
       3.0\,PD & \fpath{results/phase2\_montecarlo.mat}\\
WP2.6 & Phase 2 milestone document; manuscript revision (or follow-on TPWRD
       paper draft); D2 review. &
       8.0\,PD & revised IEEE Access response or TPWRD draft; D-C achieved\\
\midrule
\multicolumn{2}{r}{\textbf{Phase 2 total}} & \textbf{28.0\,PD} & \\
\bottomrule
\end{longtable}

\subsection{Phase 3 -- 3-phase / IEEE feeders / Taylor-Fourier (W10--24; $\approx 50$\,PD)}

\begin{longtable}{L{1.2cm}L{9.4cm}C{1.2cm}L{3.2cm}}
\toprule\rowcolor{sambplight}
\textbf{WP} & \textbf{Sub-task} & \textbf{Effort} & \textbf{Output}\\\midrule\endhead
WP3.1 & Generalise state-space to 3-port admittance $Y_{abc}(j\omega_0)$;
       multi-port closed-form $H$ per Phase 2. &
       8.0\,PD & \fpath{Y\_abc\_closedform.m}\\
WP3.2 & Add laterals, tap loads, $\geq 1$ DG; parameterise upstream Thevenin
       source. &
       5.0\,PD & extended feeder model\\
WP3.3 & Build IEEE 13- / 34- / 123-node test-feeder digital twins
       (PSCAD + MATLAB). &
       10.0\,PD & \fpath{IEEE\_13.pscx}, \fpath{IEEE\_34.pscx},
       \fpath{IEEE\_123.pscx}\\
WP3.4 & Add SLG, LL, LLG fault types in addition to single-phase HIF;
       re-run extended 720-grid. &
       5.0\,PD & \fpath{results/phase3\_faulttypes.mat}\\
WP3.5 & Replace single-bin DFT with first-order Taylor-Fourier; quantify
       phasor-bias improvement on arc-modulated waveforms; map
       $J(\alpha, R_x)$ on representative case
       (Hermann-Krener / STRIKE-GOLDD). &
       8.0\,PD & \fpath{tft\_phasor.m}; \fpath{fig/J\_surface\_*.pdf};
       identifiability map\\
WP3.6 & Multi-port FIM (3 ports $\rightarrow$ 6 real obs per port);
       project onto $(\alpha, R_x)$; update ML cost and CRLB overlays. &
       6.0\,PD & \fpath{fim\_multiport.m}\\
WP3.7 & Validate on CNRS-2024 IEEE-34 HIF dataset
       (DOI 10.57745/KRYCYY); compare vs simulator output. &
       4.0\,PD & \fpath{phase3\_cnrs\_validation.csv}\\
WP3.8 & Phase 3 milestone; conference paper (IEEE PES GM or ISGT); D3 review. &
       4.0\,PD & conference paper draft; D-D achieved\\
\midrule
\multicolumn{2}{r}{\textbf{Phase 3 total}} & \textbf{50.0\,PD} & \\
\bottomrule
\end{longtable}

\subsection{Phase 4 -- Impairments + competitor benchmark + arc diversity (W18--30; $\approx 45$\,PD)}

\begin{longtable}{L{1.2cm}L{9.4cm}C{1.2cm}L{3.2cm}}
\toprule\rowcolor{sambplight}
\textbf{WP} & \textbf{Sub-task} & \textbf{Effort} & \textbf{Output}\\\midrule\endhead
WP4.1 & Add 5 impairment classes: impulsive (Bernoulli-Gaussian), harmonic
       (2/5/7/11), CT/PT saturation per IEEE C37.110, $\pm 0.5$\,Hz off-nominal
       per IEEE C37.118.1 P-class, 12/14/16-bit ADC quantisation; re-run
       720-grid for each. &
       10.0\,PD & \fpath{impairments\_lib.m};
       \fpath{phase4\_impairments.mat}\\
WP4.2 & Cassie-Mayr-Kizilcay arc model (Darwish-Elkalashy-2005); re-run grid;
       compare estimator performance to diode arc. &
       5.0\,PD & \fpath{arc\_kizilcay.m}\\
WP4.3 & Integrate Wang-2020 distortion-controllable arc (open-source
       PSCAD-FILE-DISTC-HIAF-Model); use as PSCAD test stimulus. &
       5.0\,PD & \fpath{pscad/wang2020\_arc/}\\
WP4.4 & Torres-2022 stochastic-configurable arc (build-up, shoulder,
       asymmetry, avalanche, intermittence, modulation); re-run grid. &
       5.0\,PD & \fpath{arc\_torres.m}\\
WP4.5 & Re-implement four competitors on identical 720-case waveforms:
       (1)~Paramo-2023 eigenvalue, (2)~Iurinic-2018 / Orozco-Henao-2020
       spectral, (3)~Cui-Weng-2020 $\mu$-PMU, (4)~Zeng-2021 double-ended HIF.
       Publish numerical Table 3-bis. &
       12.0\,PD & \fpath{matlab/competitors/};
       \fpath{phase4\_table3bis.csv}\\
WP4.6 & Phase 4 milestone; benchmark-data paper (IEEE TSG candidate);
       D4 review. &
       8.0\,PD & TSG draft; D-E achieved\\
\midrule
\multicolumn{2}{r}{\textbf{Phase 4 total}} & \textbf{45.0\,PD} & \\
\bottomrule
\end{longtable}

\subsection{Phase 5 -- HIL pilot + DTaaS productisation (W26--48; $\approx 80$\,PD)}

\begin{longtable}{L{1.2cm}L{9.4cm}C{1.2cm}L{3.2cm}}
\toprule\rowcolor{sambplight}
\textbf{WP} & \textbf{Sub-task} & \textbf{Effort} & \textbf{Output}\\\midrule\endhead
WP5.1 & RTDS / Typhoon HIL platform commissioning at IIT-Madras lab;
       engagement memos with NUS GEMS, NTU CTSP, Amprion. &
       15.0\,PD & HIL access agreement; cabling diagram\\
WP5.2 & Real relay-class IED integration (SEL / ABB / Siemens) via IEC
       61850-9-2 SV; verify SV stream timing and GOOSE round-trip. &
       20.0\,PD & \fpath{IED\_integration\_log.md}; SV timing report\\
WP5.3 & Live HIF arc HIL scenarios using Wang-2020 distortion-controllable arc;
       $\geq 25$ scenarios; report end-to-end estimate latency $<5$ cycles. &
       15.0\,PD & \fpath{hil/scenario\_results.csv}\\
WP5.4 & SAMBPS DTaaS Protection-Validation module: API spec (REST + Python
       client), DT scenario-engine integration, identifiability + CRLB
       overlays in UI, documentation, commercial case study. &
       20.0\,PD & \fpath{dtaas/protection\_validation/}; API docs;
       case study\\
WP5.5 & Second journal paper (IEEE TPWRD / TSG) bringing HIL evidence +
       4-competitor benchmark + multi-port CRLB. &
       10.0\,PD & TPWRD/TSG submission; D-F, D-G, D-H achieved\\
\midrule
\multicolumn{2}{r}{\textbf{Phase 5 total}} & \textbf{80.0\,PD} & \\
\bottomrule
\end{longtable}

\paragraph{Total project effort.}
Phase 0 $\approx 14$ + Phase 1 $\approx 35$ + Phase 2 $\approx 28$ +
Phase 3 $\approx 50$ + Phase 4 $\approx 45$ + Phase 5 $\approx 80$
$=\approx 252$ person-days.  With one full-time engineer this is $\sim 50$
weeks; with two engineers in parallel from W2 onward (Phase 1 + Phase 2
overlap, and Phase 3 + Phase 4 overlap), the calendar fits the 48-week
envelope with a small contingency margin.

% ===========================================================================
\section{RACI matrix}
% ===========================================================================

\textbf{R}esponsible (does the work),
\textbf{A}ccountable (signs off),
\textbf{C}onsulted (input before decision),
\textbf{I}nformed (notified after).
\textbf{PI} = Anoop Eluvathingal (SAMBPS PI).
\textbf{LE} = Arjundas K.\ (lead engineer).
\textbf{HA} = Prof.\ K.\ Shanthi Swarup (host advisor, IIT-Madras).
\textbf{E2} = second engineer (Sanjana Verma / V.\ Pranjal / Renju PB; allocated per WP).
\textbf{ER} = Prof.\ Christian Rehtanz (TU Dortmund EMT reviewer).
\textbf{FE} = Fabian Erlemeyer (Amprion HVDC advisor).
\textbf{DS} = Prof.\ Dipti Srinivasan (NUS GEMS).
\textbf{YX} = Prof.\ Yan Xu (NTU).
\textbf{IND} = Rajkumar Palaniappan / D.\ Hilbrich (IED industry partners).

\renewcommand{\arraystretch}{1.10}
\begin{longtable}{L{6.0cm} *{9}{C{0.7cm}}}
\toprule\rowcolor{sambplight}
\textbf{Work package} & PI & LE & HA & E2 & ER & FE & DS & YX & IND\\\midrule\endhead
WP0.1 Editorial / metric harmonisation       & A & R & C & I & -- & -- & -- & -- & --\\
WP0.2 Prior-art restructure                  & A & R & C & I & -- & -- & -- & -- & --\\
WP0.3 Reference set + DOI pass               & I & R & C & R & -- & -- & -- & -- & --\\
WP0.4 Public repo + CI                       & A & R & I & R & -- & -- & -- & -- & --\\
WP0.5 Appendix A: derivation + $\partial H/\partial\theta$ & A & R & C & -- & -- & -- & -- & -- & --\\
WP0.6 IEEE Access camera-ready               & A & R & C & I & -- & -- & -- & -- & --\\
WP1.1 PSCAD model                            & A & R & C & R & C & -- & -- & -- & --\\
WP1.2 EMTP-RV mirror                         & A & C & C & R & R & -- & -- & -- & --\\
WP1.3 50-section reference state-space       & A & R & C & -- & -- & -- & -- & -- & --\\
WP1.4 Optimiser re-run on 3 simulators       & A & R & C & R & -- & -- & -- & -- & --\\
WP1.5 Monte-Carlo + 95-pct CI                & A & R & I & R & -- & -- & -- & -- & --\\
WP1.6 Proper-complex-Gaussian-ratio FIM      & A & R & C & -- & -- & -- & -- & -- & --\\
WP1.7 Phase 1 milestone + arXiv              & A & R & C & I & -- & -- & -- & -- & --\\
WP2.1 Closed-form $H$ derivation             & A & R & C & -- & C & -- & -- & -- & --\\
WP2.2 $\partial H/\partial\alpha$ closed-form & A & R & C & -- & -- & -- & -- & -- & --\\
WP2.3 50-section reproduction                & A & R & I & R & -- & -- & -- & -- & --\\
WP2.4 Analytical gradients in optimiser      & A & R & I & R & -- & -- & -- & -- & --\\
WP2.5 720 + MC re-run                        & A & R & I & R & -- & -- & -- & -- & --\\
WP2.6 Phase 2 milestone                      & A & R & C & I & -- & -- & -- & -- & --\\
WP3.1 $Y_{abc}$ closed-form                  & A & R & C & -- & C & -- & -- & -- & --\\
WP3.2 Laterals / DG                          & A & R & C & R & -- & -- & -- & -- & --\\
WP3.3 IEEE 13/34/123 feeders                 & A & R & C & R & -- & -- & -- & -- & --\\
WP3.4 Fault types                            & A & R & I & R & -- & -- & -- & -- & --\\
WP3.5 Taylor-Fourier + identifiability map   & A & R & C & -- & -- & -- & -- & -- & --\\
WP3.6 Multi-port FIM                         & A & R & C & -- & -- & -- & -- & -- & --\\
WP3.7 CNRS dataset validation                & A & R & I & R & -- & -- & -- & -- & --\\
WP3.8 Phase 3 milestone                      & A & R & C & I & I & -- & -- & -- & --\\
WP4.1 Field-grade impairments                & A & R & C & R & -- & -- & -- & -- & --\\
WP4.2 Cassie-Mayr-Kizilcay arc               & A & R & C & -- & C & -- & -- & -- & --\\
WP4.3 Wang-2020 distortion-controllable      & A & R & I & R & -- & -- & -- & -- & --\\
WP4.4 Torres-2022 stochastic arc             & A & R & I & R & -- & -- & -- & -- & --\\
WP4.5 Competitor benchmark (4 methods)       & A & R & C & R & -- & -- & -- & C & --\\
WP4.6 Phase 4 milestone                      & A & R & C & I & I & -- & -- & -- & --\\
WP5.1 HIL platform commissioning             & A & R & C & R & -- & -- & C & C & C\\
WP5.2 Real-IED integration (IEC 61850-9-2)   & A & R & C & -- & -- & -- & -- & -- & R\\
WP5.3 Live HIF arc HIL                       & A & R & C & R & -- & C & C & C & C\\
WP5.4 SAMBPS DTaaS module v1.0               & A & R & I & R & -- & -- & I & I & I\\
WP5.5 Second journal paper                   & A & R & C & I & C & C & C & C & I\\
\bottomrule
\end{longtable}

% ===========================================================================
\section{Acceptance test specifications}
% ===========================================================================

Each deliverable D-A through D-H is locked to a quantitative acceptance
test that CI executes on every commit.  A deliverable is not signed off
until its test passes with no exceptions.

\renewcommand{\arraystretch}{1.18}
\begin{longtable}{L{2.2cm}L{1.0cm}L{2.8cm}L{5.6cm}L{2.8cm}}
\toprule\rowcolor{sambplight}
\textbf{Deliverable} & \textbf{Test ID} & \textbf{Input dataset} &
\textbf{Acceptance predicate} & \textbf{CI artefact}\\\midrule\endhead

D-A IEEE Access v2.0 & T-A1 &
Camera-ready PDF + supplementary &
(a) Single coherent metric set; (b) Appendix A includes
$\pi$-derivation + $\partial H/\partial\theta$;
(c) repository link resolves; (d) references $\geq 35$ entries with
valid DOIs. &
\fpath{tests/checklist\_camera\_ready.py}\\

D-B Phase-1 cross-validation & T-B1 &
720-grid on PSCAD + EMTP + 50-sec ref &
Mean location error $<2$\,\% at $\mathrm{SNR}_I \geq 30$\,dB across
all 3 simulators; max $<5$\,\%; per-cell 95\,\% CI excludes zero in
$<5$\,\% of cells. &
\fpath{test\_phase1\_crossplatform.py}\\

D-B (CRLB) & T-B2 &
Synthetic Monte-Carlo @ SNR $\in\{20,30,40\}$ dB &
Empirical RMS error within 25\,\% of proper-complex-Gaussian-ratio CRLB at
SNR $\geq 30$\,dB; Geary-Hinkley validity test reported per cell. &
\fpath{test\_crlb\_proper.py}\\

D-C Continuously parametrised optimiser & T-C1 &
Synthetic 720-grid &
Modelling error vs 50-section ref $<5$\,\% across all $(\alpha, R_x)$;
estimator improvement $\geq 30$\,\% at $\mathrm{SNR}_I \leq 30$\,dB
vs Phase-1 baseline. &
\fpath{test\_phase2\_modelfit.py}\\

D-D Three-phase + IEEE feeders + TFT & T-D1 &
IEEE 13/34/123-node + CNRS dataset &
Mean location error $<3$\,\% on IEEE 34-node at SNR $\geq 30$\,dB;
phasor-bias improvement vs single-bin DFT $\geq 50$\,\% on
arc-modulated waveform; identifiability map published. &
\fpath{test\_phase3\_3phase.py}\\

D-E Impairments + competitor benchmark & T-E1 &
720-grid $\times$ 5 impairment classes &
$<5$\,\% mean location error at SNR $\geq 30$\,dB across all 5
impairment classes; numerical superiority over $\geq 2$ of 4
competitors. &
\fpath{test\_phase4\_impairments.py},\newline
\fpath{test\_phase4\_benchmark.py}\\

D-F HIL pilot & T-F1 &
$\geq 25$ live HIF scenarios on RTDS / Typhoon HIL &
Real-IED end-to-end estimate latency $<5$ power-frequency cycles;
field-style HIF arc reproduced; $\geq 1$ institutional pilot signoff. &
\fpath{hil/test\_latency.py};\newline pilot signoff letter\\

D-G SAMBPS DTaaS module v1.0 & T-G1 &
Smoke test on three reference twins &
Module loads in DT scenario engine; API endpoints respond;
identifiability + CRLB overlays render; smoke test passes on
Reference Twins 1--3 (11 kV / 100 km, IEEE 34-node, HVDC stub). &
\fpath{dtaas/tests/smoke\_test.py}\\

D-H Second journal paper & T-H1 &
Submission record &
Submitted to IEEE TPWRD or TSG; reproducibility repo cited; HIL
evidence included; $\geq 4$ competitors benchmarked numerically. &
Submission ID;
\fpath{docs/journal\_v2\_submission.md}\\

\bottomrule
\end{longtable}

% ===========================================================================
\section{Schedule}
% ===========================================================================

\subsection{Phase 0 day-level schedule (W0 = 11 May 2026)}

\renewcommand{\arraystretch}{1.10}
\begin{longtable}{L{1.4cm}L{1.5cm}L{7.8cm}L{1.2cm}L{2.6cm}}
\toprule\rowcolor{sambplight}
\textbf{Day} & \textbf{Date} & \textbf{Lead activity} & \textbf{Owner} & \textbf{Output}\\\midrule\endhead
W0 Mon & 11 May & Title + abstract rewrite (WP0.1)                              & LE     & Draft v0\\
W0 Tue & 12 May & Conclusion roadmap + metric harmonisation (WP0.1)             & LE     & Draft v1\\
W0 Wed & 13 May & Prior-art restructure into 7 families (WP0.2)                 & LE     & Revised \S\,I\\
W0 Thu & 14 May & Wildfire / safety motivation paragraph + forward-reference   & LE     & Revised \S\,I.1\\
W0 Fri & 15 May & Reference set expansion start: pull v2 \S\,7 bibliography to BibTeX (WP0.3) & E2 & Initial \texttt{.bib}\\
W1 Mon & 18 May & Reference DOI + IEEE-style consistency pass (WP0.3)           & E2     & Cleaned \texttt{.bib}\\
W1 Tue & 19 May & Public repo standup: README, LICENSE, .gitignore, CI skeleton (WP0.4) & LE & Repo (private)\\
W1 Wed & 20 May & Push existing .m source + 720-case MAT + figure scripts (WP0.4) & LE   & First commit\\
W1 Thu & 21 May & Symbolic \texttt{sym/eig} regression test for $A$; \texttt{tic/toc} CPU report (WP0.4 sub-task 8) & LE & Test passes\\
W1 Fri & 22 May & Appendix A start: annotated $\pi$-circuit + KVL/KCL (WP0.5)   & LE     & AppendixA v0\\
W2 Mon & 25 May & Symbolic $\partial H/\partial\alpha$ and $\partial H/\partial R_x$ (WP0.5) & LE & \fpath{derive\_partials.m}\\
W2 Tue & 26 May & Arc-parameter provenance citations + numeric line params with units (WP0.5) & LE & AppendixA v1\\
W2 Wed & 27 May & Co-author signoff round (PI + HA review)                      & PI/HA  & Comments\\
W2 Thu & 28 May & Final IEEE Access camera-ready integration; submit (WP0.6)   & LE/PI  & Submission ID\\
W2 Fri & 29 May & \textbf{D0 review meeting}; flip repo to public; mint Zenodo DOI; kick off Phase 1 & PI/LE/HA & D-A signed off\\
\bottomrule
\end{longtable}

\subsection{Phases 1--5 week-level schedule}

\renewcommand{\arraystretch}{1.10}
\begin{longtable}{L{1.4cm}C{1.1cm}L{8.0cm}L{3.8cm}}
\toprule\rowcolor{sambplight}
\textbf{Week} & \textbf{Phase} & \textbf{Active WPs} & \textbf{Milestone}\\\midrule\endhead
W2--4   & 1   & WP1.1 PSCAD model                                               & --\\
W4--6   & 1   & WP1.2 EMTP-RV mirror; WP1.3 50-section reference                & --\\
W5--7   & 1   & WP1.4 cross-platform optimiser re-run                           & --\\
W6--9   & 1+2 & WP1.5 Monte-Carlo; WP1.6 proper-complex-Gaussian-ratio FIM (parallel); WP2.1 closed-form $H$ starts & --\\
W9--10  & 1   & WP1.7 Phase 1 milestone; arXiv preprint                         & \textbf{D1 / D-B}\\
W10--12 & 2   & WP2.2 $\partial H/\partial\alpha$ closed-form; WP2.3 50-section reproduction & --\\
W12--14 & 2   & WP2.4 analytical gradients; WP2.5 720 + MC re-run               & --\\
W15--16 & 2   & WP2.6 Phase 2 milestone                                         & \textbf{D2 / D-C}\\
W10--14 & 3   & WP3.1 $Y_{abc}$; WP3.2 laterals / DG; WP3.3 IEEE 13/34/123 (rolling) & --\\
W14--18 & 3   & WP3.4 fault types; WP3.5 Taylor-Fourier; WP3.6 multi-port FIM   & --\\
W18--22 & 3   & WP3.7 CNRS validation                                           & --\\
W22--24 & 3   & WP3.8 Phase 3 milestone; conference paper                       & \textbf{D3 / D-D}\\
W18--20 & 4   & WP4.1 impairment classes (parallel with Phase 3 tail)           & --\\
W20--24 & 4   & WP4.2--4.4 arc-class diversification                            & --\\
W24--28 & 4   & WP4.5 competitor benchmark (4 methods)                          & --\\
W28--30 & 4   & WP4.6 Phase 4 milestone; benchmark paper                        & \textbf{D4 / D-E}\\
W26--32 & 5   & WP5.1 HIL platform commissioning                                & --\\
W32--38 & 5   & WP5.2 real-IED integration                                      & --\\
W38--42 & 5   & WP5.3 live HIF HIL scenarios                                    & --\\
W42--46 & 5   & WP5.4 SAMBPS DTaaS module v1.0                                  & --\\
W46--48 & 5   & WP5.5 second journal paper; institutional pilot signoff         & \textbf{D5 / D-F, D-G, D-H}\\
\bottomrule
\end{longtable}

% ===========================================================================
\section{Communication plan}
% ===========================================================================

\renewcommand{\arraystretch}{1.18}
\begin{longtable}{L{2.3cm}L{3.2cm}L{3.6cm}L{6.0cm}}
\toprule\rowcolor{sambplight}
\textbf{Cadence} & \textbf{Forum} & \textbf{Participants} & \textbf{Standing agenda}\\\midrule\endhead
Daily (15 min)         & Stand-up (Slack / Zoom) & LE + E2 &
What's done since yesterday; what's planned today; blockers.\\
Weekly (60 min)        & Project sync            & PI + LE + HA + E2 &
WP burn-down; CI status; risk register review; upcoming-week plan.\\
Bi-weekly (45 min)     & Technical review        & LE + HA + ER (Phase 1/4) &
Walk through new code, derivations, results; HA / ER signoff or
revision request.\\
Phase-end (90 min)     & Decision-gate review    & PI + LE + HA + relevant
external partners &
Acceptance-test results vs predicates; signoff or return to prior WP;
publication plan.\\
Quarterly (90 min)     & SAMBPS leadership update & PI + family co-founders + advisors &
Phase progress; budget; productisation pathway; partnership pipeline.\\
Ad-hoc                 & Escalation              & PI + HA &
Triggered by R1, R2, R6, R9, R10 (per v2 risk register) or any blocker
exceeding 2 working days.\\
\bottomrule
\end{longtable}

\subsection{Decision-gate (D0--D5) review template}

\begin{enumerate}
  \item \textbf{Phase summary} (1 paragraph). Goal of the phase, what was
    delivered.
  \item \textbf{Acceptance-test results.} One line per test ID
    (T-A1, T-B1, \ldots) with PASS / FAIL and the measured value vs the
    predicate.
  \item \textbf{Risk-register update.} Any v2 risk that changed likelihood
    or impact this phase; new risks added.
  \item \textbf{KPI snapshot.} Current value of every v2 KPI vs the SAMBPS
    DTaaS v1.0 target.
  \item \textbf{Decision.} One of:
    (a)~\emph{Approve progression} to next phase;
    (b)~\emph{Conditional approval} with explicit follow-on items and dates;
    (c)~\emph{Return} to a prior WP, with re-entry plan.
  \item \textbf{Publication artefact.} Link to the external artefact
    produced by the phase (camera-ready, arXiv preprint, conference paper,
    journal submission, repo release).
  \item \textbf{Signoff.} PI and HA signatures + date.
\end{enumerate}

% ===========================================================================
\section{Repository and artefact inventory}
% ===========================================================================

The single source of truth for this sub-project is the new folder
\texttt{04\_code/sambp/fault\_location\_id/}, created alongside this plan
and structured to mirror the four existing SAMBP sub-projects
(\texttt{sync\_oc}, \texttt{transformer\_diff}, \texttt{line\_diff},
\texttt{bus\_diff}).  The public mirror at
\texttt{github.com/SAMBPS-DTaaS/HIF-TF-Locator} is minted from this folder
at decision gate D0.

\begin{verbatim}
04_code/sambp/fault_location_id/
|-- README.md                                # project overview, citation
|-- __init__.py
|-- docs/
|   |-- FaultLocationIdentification_ExecutionPlan.tex
|   |-- FaultLocationIdentification_ExecutionPlan.pdf   # this document
|   |-- AppendixA_derivation.tex             # WP0.5
|   |-- AppendixB_correctedCRLB.tex          # WP1.6
|   |-- changelog.md
|-- models/
|   |-- __init__.py
|   |-- faultloc_pi_section_model.py         # WP0.5  (Phase 0)
|   |-- faultloc_distributed_param_model.py  # WP2.1  (Phase 2)
|   |-- faultloc_three_phase_model.py        # WP3.1  (Phase 3)
|   |-- faultloc_ieee_feeders.py             # WP3.3  (Phase 3)
|   |-- faultloc_taylor_fourier.py           # WP3.5  (Phase 3)
|   |-- faultloc_arc_models.py               # WP4.2-4.4 (Phase 4)
|   `-- faultloc_noise_impairments.py        # WP4.1  (Phase 4)
|-- inverse_estimation/
|   |-- __init__.py
|   |-- faultloc_two_stage_optimiser.py      # WP2.4
|   |-- faultloc_analytical_gradients.py     # WP2.2
|   |-- faultloc_crlb_proper.py              # WP1.6
|   |-- faultloc_crlb_dualchannel.py         # WP1.6
|   `-- faultloc_fim_multiport.py            # WP3.6
|-- adaptation/
|   |-- __init__.py
|   |-- faultloc_identifiability_check.py    # WP3.5
|   `-- faultloc_confidence_gate.py          # Phase 5 hand-off to DTaaS
|-- evaluation/
|   |-- __init__.py
|   |-- faultloc_competitor_paramo.py        # WP4.5
|   |-- faultloc_competitor_iurinic.py       # WP4.5
|   |-- faultloc_competitor_cuiweng.py       # WP4.5
|   `-- faultloc_competitor_zeng.py          # WP4.5
|-- outputs/                                 # CSV / MAT / PNG artefacts
|-- run_faultloc_phase0_baseline.py
|-- run_faultloc_phase1_crossplatform.py
|-- run_faultloc_phase2_continuous_param.py
|-- run_faultloc_phase3_threephase.py
|-- run_faultloc_phase4_impairments.py
`-- run_faultloc_phase5_hil.py
\end{verbatim}

\paragraph{Naming convention.}
\texttt{phase\{N\}\_\{topic\}.\{ext\}} for results;
\texttt{faultloc\_*.py} for code, mirroring the
\texttt{line\_*.py} / \texttt{transformer\_*.py} / \texttt{bus\_*.py}
prefixes used by the other SAMBP sub-projects.  Every commit references
at least one WP ID in the message; CI enforces this.  Public flip occurs
at D0 (W2 Fri); Zenodo DOIs are minted on every D0 / D2 / D4 / D5
release.  The Data Availability Statement in the IEEE Access camera-ready
and the second journal paper points to the Zenodo DOI plus the GitHub
head.

% ===========================================================================
\section{Risk register}
% ===========================================================================

\renewcommand{\arraystretch}{1.10}
\begin{longtable}{L{0.6cm}L{4.2cm}L{1.2cm}L{1.2cm}L{6.4cm}}
\toprule\rowcolor{sambplight}
\textbf{ID} & \textbf{Risk} & \textbf{Likely} & \textbf{Impact} & \textbf{Mitigation / closing WP}\\\midrule\endhead
R1 & Gaussian-on-$H$ FIM is invalid in the HIF regime
   & High & High
   & WP1.6 re-derives FIM via proper-complex-Gaussian-ratio +
     joint dual-channel FIM; CRLB-vs-empirical overlays.\\
R2 & 39.44\,\% section-modelling-error ceiling caps estimator accuracy
   & High & High
   & WP2.1--2.5 closed-form distributed-parameter $H$; WP2.4 analytical
     gradients in optimiser.\\
R3 & Single-section / single-phase model mismatches real feeders
   & Med  & High
   & WP3.1--3.4 generalise to 3-phase $Y_{abc}$ + IEEE 13/34/123 + SLG /
     LL / LLG fault types.\\
R4 & Diode arc parameter provenance gap
   & Med  & Med
   & WP0.5 cite source of every parameter; WP4.2--4.4 add 3 alternative
     arc classes; report $\Delta$-error.\\
R5 & Single-bin DFT biased on non-stationary HIF
   & High & Med
   & WP3.5 first-order Taylor-Fourier estimator; identifiability map.\\
R6 & Categorical comparison with prior-art is non-quantitative
   & High & Med
   & WP4.5 numerical Table 3-bis with four competitors on identical
     720-case waveforms.\\
R7 & Symbolic derivation errors propagate to optimiser
   & Low  & High
   & WP0.4 \texttt{sym/eig} regression test in CI; WP0.5 verify in MATLAB
     \texttt{sym/diff}.\\
R8 & HIL access at IIT-Madras delayed by external schedule
   & Med  & High
   & WP5.1 engagement memos with NUS GEMS / NTU CTSP / Amprion as
     redundant access paths; Typhoon HIL accepted as RTDS substitute.\\
R9 & Geary-Hinkley validity test fails in part of the operating range
   & Med  & High
   & WP1.6 publish per-cell validity flag; flag down-stream usage; fall
     back to joint dual-channel FIM where ratio is unstable.\\
R10 & Real-world HIF is more random than diode arc captures
   & Med  & Med
   & WP4.3 / WP4.4 / WP5.3 Wang-2020 + Torres-2022 + live HIL arc.\\
\bottomrule
\end{longtable}

% ===========================================================================
\section{KPI dashboard}
% ===========================================================================

The 17 KPIs imported from the v2 plan, with the Phase that drives them
to the SAMBPS DTaaS v1.0 target.

\renewcommand{\arraystretch}{1.18}
\begin{longtable}{L{0.7cm}L{6.6cm}L{2.4cm}L{2.4cm}L{1.4cm}}
\toprule\rowcolor{sambplight}
\textbf{\#} & \textbf{KPI} & \textbf{Baseline} & \textbf{Target} & \textbf{Phase}\\\midrule\endhead
K01 & Mean location error, AWGN, $\mathrm{SNR}_I=30$ dB & 1.18\,\%             & $<1.0$\,\%  & 1\\
K02 & Max location error, AWGN, $\mathrm{SNR}_I=30$ dB  & --                   & $<5$\,\%    & 1\\
K03 & Modelling error vs 50-section reference           & 39.44\,\%            & $<5$\,\%    & 2\\
K04 & Estimator improvement at $\mathrm{SNR}_I\le 30$ dB & --                  & $\geq 30$\,\% & 2\\
K05 & Mean location error, IEEE 34-node                 & --                   & $<3$\,\%    & 3\\
K06 & Phasor-bias improvement vs single-bin DFT         & --                   & $\geq 50$\,\% & 3\\
K07 & Mean location error across 5 impairment classes   & --                   & $<5$\,\%    & 4\\
K08 & \# competitors numerically beaten on location error & 0                  & $\geq 2$ of 4 & 4\\
K09 & Real-IED end-to-end estimate latency              & --                   & $<5$ cycles & 5\\
K10 & \# institutional HIL pilot signoffs               & 0                    & $\geq 1$    & 5\\
K11 & References, IEEE-style, valid DOI                 & 19 unique            & $\geq 35$   & 0\\
K12 & Public repo + Zenodo DOI                          & none                 & live at D0  & 0\\
K13 & CI green build rate (rolling 30 days)             & --                   & $\geq 95$\,\% & 0+\\
K14 & 95th-percentile estimator CPU time                & not reported         & published   & 0\\
K15 & Hyperparameter sensitivity table coverage         & not present          & full        & 0\\
K16 & DTaaS Protection-Validation module smoke test     & not built            & passes      & 5\\
K17 & Second journal paper submitted                    & 0                    & 1           & 5\\
\bottomrule
\end{longtable}

% ===========================================================================
\section{Bottom line and approval block}
% ===========================================================================

\noindent\colorbox{sambplight}{\parbox{\dimexpr\textwidth-2\fboxsep\relax}{%
\vspace{0.1cm}\color{sambpblue}
The IEEE Access manuscript is publishable today after Phase 0 ($\approx 14$\,PD
over W0--2).  All four shortcomings flagged in \S\,3 close inside Phase 0
(D-A).  The two structural risks identified by the v2 plan close in Phase 1
(corrected CRLB $\rightarrow$ D-B) and Phase 2 (continuously parametrised
optimiser $\rightarrow$ D-C).  Phases 3--5 then convert the manuscript into
a literature-anchored, field-validated, three-phase, hardware-tested SAMBPS
DTaaS Protection-Validation module with a second journal paper (D-H) and
an institutional HIL pilot (D-F).  Total horizon 48 weeks; total effort
$\approx 252$\,PD; staffing assumption is one full-time lead engineer + one
half-time second engineer + advisor time per the RACI in \S\,5.
\vspace{0.1cm}
}}

\vspace{0.4cm}

\subsection*{Approval block}

\renewcommand{\arraystretch}{1.6}
\begin{longtable}{L{5.0cm}L{4.6cm}L{3.0cm}L{2.0cm}}
\toprule\rowcolor{sambplight}
\textbf{Role} & \textbf{Name} & \textbf{Signature} & \textbf{Date}\\\midrule
Principal Investigator                  & Anoop Eluvathingal                   & & \\
Lead Engineer                           & Arjundas K.                          & & \\
Host Advisor                            & Prof.\ K.\ Shanthi Swarup, IIT Madras & & \\
EMT Cross-Validation Reviewer           & Prof.\ Christian Rehtanz, TU Dortmund & & \\
HVDC Application Advisor                & Fabian Erlemeyer, Amprion            & & \\
Singapore Academic Anchor (smart grid)  & Prof.\ Dipti Srinivasan, NUS GEMS    & & \\
Singapore Academic Anchor (stability)   & Prof.\ Yan Xu, NTU                   & & \\
\bottomrule
\end{longtable}

\vspace{1.0cm}
\hrule
\vspace{0.2cm}
\centerline{\itshape\small -- End of Fault Location Identification Project Execution Plan --}

\end{document}
